\documentclass[aspectratio=169]{beamer}
\usetheme{Madrid}
\usecolortheme{default}
\usepackage[utf8]{inputenc}
\usepackage[T5]{fontenc}
\usepackage{graphicx}
\usepackage{hyperref}
\usepackage{booktabs}
\usepackage{amsmath}

\setbeamertemplate{caption}[numbered]
\renewcommand{\figurename}{Hình}
\renewcommand{\thefigure}{11.\arabic{figure}}

\title[Chương 11: Phân vùng Hình ảnh]{Học Sâu (Deep Learning) \\ \vspace{0.5cm} \Large Chương 11: Phân vùng Hình ảnh \\ (Image Segmentation)}
\author{Đại học Đông Á}
\date{\today}

\begin{document}

\begin{frame}
    \titlepage
\end{frame}

\begin{frame}{Nội dung chương}
    \tableofcontents
\end{frame}

\section{1. Tổng quan các Bài toán Thị giác Máy tính}

\begin{frame}{Ba nhánh chính của Thị giác Máy tính}
    \begin{columns}
        \column{0.4\textwidth}
        \begin{itemize}
            \item \textbf{Image Classification (Phân loại):} Gán nhãn cho toàn bộ bức ảnh ("Ảnh này có con mèo").
            \item \textbf{Object Detection (Phát hiện vật thể):} Vẽ hộp giới hạn (Bounding box) quanh vật thể và phân loại chúng ("Con mèo ở tọa độ x, y").
            \item \textbf{Image Segmentation (Phân vùng):} Phân loại ở cấp độ \textbf{Điểm ảnh (Pixel)}. Cắt chính xác đường viền của vật thể ("Những pixel nào thuộc về con mèo?").
        \end{itemize}
        
        \column{0.6\textwidth}
        \begin{figure}
            \centering
            \includegraphics[width=\textwidth,height=0.6\textheight,keepaspectratio]{../../images/ch11/computer_vision_tasks.da2bf0ea.png}
            \caption{Ba bài toán cốt lõi trong CV}
        \end{figure}
    \end{columns}
\end{frame}

\begin{frame}{Phân loại Bài toán Image Segmentation}
    \begin{columns}
        \column{0.5\textwidth}
        \begin{itemize}
            \item \textbf{Semantic Segmentation:} Phân loại pixel theo "Danh mục" (Category). Nếu có 2 con mèo, cả 2 sẽ có cùng màu (cùng nhãn "Mèo").
            \item \textbf{Instance Segmentation:} Phân loại pixel theo "Cá thể" (Instance). Phân biệt rạch ròi "Con mèo 1" và "Con mèo 2".
            \item \textbf{Panoptic Segmentation:} Sự kết hợp cao cấp của cả Semantic và Instance Segmentation.
        \end{itemize}
        
        \column{0.5\textwidth}
        \begin{figure}
            \centering
            \includegraphics[width=\textwidth,height=0.6\textheight,keepaspectratio]{../../images/ch11/instance_segmentation.818c62ba.png}
            \caption{Semantic vs Instance Segmentation}
        \end{figure}
    \end{columns}
\end{frame}

\section{2. Xây dựng Mô hình Phân vùng Semantic}

\begin{frame}{Bài toán: Phân tách Tiền cảnh và Hậu cảnh}
    \begin{columns}
        \column{0.5\textwidth}
        \begin{itemize}
            \item Ta sẽ xây dựng một mô hình CNN từ đầu để thực hiện Semantic Segmentation trên tập dữ liệu Oxford-IIIT Pets (Chó \& Mèo).
            \item \textbf{Mục tiêu:} Phân biệt đâu là chủ thể (Vật nuôi), đâu là đường viền, và đâu là hậu cảnh (Background).
            \item Ảnh đầu vào: $200 \times 200 \times 3$.
        \end{itemize}
        
        \column{0.5\textwidth}
        \begin{figure}
            \centering
            \includegraphics[width=\textwidth,height=0.6\textheight,keepaspectratio]{../../images/ch11/segmentation_input.d246cf5a.png}
            \caption{Ảnh đầu vào}
        \end{figure}
    \end{columns}
\end{frame}

\begin{frame}{Khái niệm Mặt nạ (Mask)}
    \begin{columns}
        \column{0.5\textwidth}
        \begin{itemize}
            \item Dữ liệu nhãn (Ground Truth) trong bài toán này không phải là một con số, mà là một \textbf{Mặt nạ (Mask)}.
            \item Kích thước Mask tương đương ảnh gốc ($200 \times 200 \times 1$).
            \item Mỗi pixel mang 1 trong 3 giá trị:
            \begin{enumerate}
                \item Foreground (Tiền cảnh)
                \item Background (Hậu cảnh)
                \item Contour (Đường viền)
            \end{enumerate}
        \end{itemize}
        
        \column{0.5\textwidth}
        \begin{figure}
            \centering
            \includegraphics[width=\textwidth,height=0.6\textheight,keepaspectratio]{../../images/ch11/segmentation_mask.cc320651.png}
            \caption{Mặt nạ phân vùng (Ground truth)}
        \end{figure}
    \end{columns}
\end{frame}

\begin{frame}{Kiến trúc Encoder - Decoder (U-Net style)}
    Mạng nơ-ron phân vùng có hình phễu, chia làm 2 nửa:
    \begin{itemize}
        \item \textbf{Encoder (Mã hóa):} Dùng \texttt{Conv2D} với \texttt{strides=2} để \textbf{Hạ độ phân giải} ảnh xuống (Ví dụ: $200 \rightarrow 100 \rightarrow 50$). Nhằm trích xuất thông tin khái niệm (What).
        \item \textbf{Decoder (Giải mã):} Dùng \texttt{Conv2DTranspose} với \texttt{strides=2} để \textbf{Tăng độ phân giải} ảnh lên lại (Ví dụ: $50 \rightarrow 100 \rightarrow 200$). Nhằm định vị lại không gian gốc (Where).
    \end{itemize}
\end{frame}

\begin{frame}{Toán học của Transposed Convolution}
    \begin{itemize}
        \item Để tăng độ phân giải (Upsampling) mà vẫn duy trì khả năng học hỏi (Learnable), ta dùng \textbf{Tích chập Chuyển vị (Transposed Convolution)}.
        \item Nó hoạt động ngược lại với Tích chập thông thường:
        \begin{itemize}
            \item Nhận 1 giá trị đầu vào (Pixel).
            \item Nhân giá trị đó với một Bộ lọc (Kernel).
            \item Phát tán (Broadcast) kết quả ra một vùng không gian lớn hơn ở đầu ra.
        \end{itemize}
        \item Nhờ đó, mạng có thể học cách "điền" thông tin chi tiết vào những vùng ảnh đã bị thu nhỏ trước đó.
    \end{itemize}
\end{frame}

\begin{frame}{Biểu đồ Loss trong Huấn luyện}
    \begin{columns}
        \column{0.5\textwidth}
        \begin{itemize}
            \item Vì đây là phân loại 3 nhãn cho từng pixel, ta dùng Hàm mất mát: \texttt{sparse\_categorical\_crossentropy}.
            \item Mô hình hội tụ rất tốt ở giai đoạn đầu.
            \item Bắt đầu xuất hiện Overfitting sau epoch thứ 25.
        \end{itemize}
        
        \column{0.5\textwidth}
        \begin{figure}
            \centering
            \includegraphics[width=\textwidth,height=0.6\textheight,keepaspectratio]{../../images/ch11/segmentation_loss.489fa0c8.png}
            \caption{Quá trình Training Loss}
        \end{figure}
    \end{columns}
\end{frame}

\begin{frame}{Kết quả Dự đoán (Prediction)}
    \begin{columns}
        \column{0.5\textwidth}
        \begin{itemize}
            \item Trực quan hóa kết quả dự đoán của mô hình tự thiết kế.
            \item Mô hình bắt được hình dáng của con mèo khá tốt (Khu vực Tiền cảnh).
            \item Tuy nhiên, các chi tiết nhỏ ở đường viền (như tai, râu) vẫn còn nhiễu và bị mờ (Artifacts).
            \item Để đạt độ chính xác SOTA, ta cần các mô hình lớn hơn rất nhiều.
        \end{itemize}
        
        \column{0.5\textwidth}
        \begin{figure}
            \centering
            \includegraphics[width=\textwidth,height=0.6\textheight,keepaspectratio]{../../images/ch11/segmentation_test.ece7f638.png}
            \caption{Mặt nạ do mô hình dự đoán}
        \end{figure}
    \end{columns}
\end{frame}

\section{3. Siêu mô hình Nền tảng: Segment Anything (SAM)}

\begin{frame}{Kỷ nguyên của Foundation Models trong CV}
    \begin{itemize}
        \item Natural Language Processing (NLP) đã thay đổi mãi mãi với sự xuất hiện của các Foundation Models (Như GPT-3, GPT-4) - Một mô hình huấn luyện một lần, giải quyết được mọi tác vụ (Zero-shot).
        \item Vào tháng 4/2023, Meta AI công bố \textbf{Segment Anything Model (SAM)}, đánh dấu cột mốc Foundation Model đầu tiên của ngành Thị giác Máy tính (Computer Vision).
        \item SAM có khả năng phân vùng "Bất cứ thứ gì" trong ảnh, mà không cần huấn luyện lại.
    \end{itemize}
\end{frame}

\begin{frame}{Kiến trúc của SAM}
    \begin{columns}
        \column{0.5\textwidth}
        SAM gồm 3 thành phần cốt lõi:
        \begin{enumerate}
            \item \textbf{Image Encoder:} Vision Transformer (ViT) kích thước cực lớn, giải nén toàn bộ thông tin ảnh.
            \item \textbf{Prompt Encoder:} Biểu diễn các Lệnh nhắc của người dùng (Điểm, Hộp, Text).
            \item \textbf{Mask Decoder:} Một mạng nhẹ hợp nhất thông tin từ Image và Prompt để xuất ra Mặt nạ (Mask).
        \end{enumerate}
        
        \column{0.5\textwidth}
        \begin{figure}
            \centering
            \includegraphics[width=\textwidth,height=0.5\textheight,keepaspectratio]{../../images/ch11/sam_architecture.dad9dae6.png}
            \caption{Kiến trúc tổng thể của Segment Anything}
        \end{figure}
    \end{columns}
\end{frame}

\begin{frame}{Siêu Dữ liệu SA-1B}
    \begin{columns}
        \column{0.5\textwidth}
        \begin{itemize}
            \item Sức mạnh của SAM đến từ Dữ liệu. Meta AI đã xây dựng tập dữ liệu \textbf{SA-1B (Segment Anything 1-Billion)}.
            \item Khối lượng khổng lồ: \textbf{11 triệu hình ảnh} với hơn \textbf{1.1 Tỷ mặt nạ phân vùng} chất lượng cao (gấp 400 lần các tập dữ liệu trước đó).
            \item Nhờ vậy, SAM có khả năng phân vùng Zero-shot với mọi thể loại vật thể từ sinh học, phương tiện, đến y tế.
        \end{itemize}
        
        \column{0.5\textwidth}
        \begin{figure}
            \centering
            \includegraphics[width=\textwidth,height=0.6\textheight,keepaspectratio]{../../images/ch11/sa1b_example.6701768b.jpg}
            \caption{Mật độ mặt nạ trong tập dữ liệu SA-1B}
        \end{figure}
    \end{columns}
\end{frame}

\section{4. Thực hành Kỹ nghệ Prompting trên SAM}

\begin{frame}{Môi trường Thử nghiệm}
    \begin{columns}
        \column{0.5\textwidth}
        \begin{itemize}
            \item Ta sẽ thử nghiệm sức mạnh của SAM trên một bức ảnh thực tế chứa rất nhiều loại Trái cây (Chuối, Đào, Xoài).
            \item Bức ảnh này có cấu trúc phức tạp với vật thể chồng chéo và bóng đổ.
            \item Mục tiêu: Gọi API của SAM để phân vùng chúng thông qua các kiểu \textbf{Prompting (Lệnh nhắc)} khác nhau.
        \end{itemize}
        
        \column{0.5\textwidth}
        \begin{figure}
            \centering
            \includegraphics[width=\textwidth,height=0.6\textheight,keepaspectratio]{../../images/ch11/fruits.8cef44dc.png}
            \caption{Ảnh Trái cây đầu vào}
        \end{figure}
    \end{columns}
\end{frame}

\begin{frame}{Automatic Mask Generation (Không cần Prompt)}
    \begin{columns}
        \column{0.5\textwidth}
        \begin{itemize}
            \item Tính năng mạnh mẽ nhất của SAM là \textbf{Tạo Mặt nạ Tự động}.
            \item SAM sẽ rải một mạng lưới (Grid) các điểm lên toàn bộ ảnh và tiến hành phân vùng mọi vật thể có thể nhận diện được.
            \item Kết quả: Nó phân mảnh chính xác từng nải chuối, từng rổ, và nền ảnh một cách hoàn hảo (Instance Segmentation).
        \end{itemize}
        
        \column{0.5\textwidth}
        \begin{figure}
            \centering
            \includegraphics[width=\textwidth,height=0.6\textheight,keepaspectratio]{../../images/ch11/bananas_all_masks.c922b7a6.png}
            \caption{Phân vùng Tự động toàn bộ ảnh}
        \end{figure}
    \end{columns}
\end{frame}

\begin{frame}{Nhắc bằng Hộp giới hạn (Box Prompting)}
    \begin{columns}
        \column{0.5\textwidth}
        \begin{figure}
            \centering
            \includegraphics[width=\textwidth,height=0.5\textheight,keepaspectratio]{../../images/ch11/mango_box.45e1bae1.png}
            \caption{Người dùng vẽ Box màu xanh lá quanh Xoài}
        \end{figure}
        \column{0.5\textwidth}
        \begin{figure}
            \centering
            \includegraphics[width=\textwidth,height=0.5\textheight,keepaspectratio]{../../images/ch11/mango_segmented.2dfb0dae.png}
            \caption{SAM lập tức trả về Mặt nạ Xoài chính xác}
        \end{figure}
    \end{columns}
\end{frame}

\begin{frame}{Nhắc bằng Điểm đơn (Single Point Prompting)}
    \begin{columns}
        \column{0.5\textwidth}
        \begin{figure}
            \centering
            \includegraphics[width=\textwidth,height=0.5\textheight,keepaspectratio]{../../images/ch11/peach_point.432d548a.png}
            \caption{Người dùng chỉ một Tọa độ (*) vào quả Đào}
        \end{figure}
        \column{0.5\textwidth}
        \begin{figure}
            \centering
            \includegraphics[width=\textwidth,height=0.5\textheight,keepaspectratio]{../../images/ch11/peach_segmented.333556ff.png}
            \caption{SAM lan truyền điểm đó thành Mặt nạ Đào}
        \end{figure}
    \end{columns}
\end{frame}

\begin{frame}{Nhắc bằng Điểm đơn trên Quả Chuối}
    \begin{columns}
        \column{0.5\textwidth}
        \begin{itemize}
            \item Tương tự, nếu ta truyền Tọa độ điểm của một quả Chuối bị khuất.
            \item Kể cả khi quả chuối bị che bởi bóng đổ hay trái cây khác, công cụ ViT khổng lồ của SAM vẫn nội suy (Interpolate) ra đường viền hoàn mỹ.
            \item Quá trình này diễn ra theo thời gian thực (Real-time) trong trình duyệt nhờ Mask Decoder siêu nhẹ.
        \end{itemize}
        
        \column{0.5\textwidth}
        \begin{figure}
            \centering
            \includegraphics[width=\textwidth,height=0.6\textheight,keepaspectratio]{../../images/ch11/banana_segmented.8e0b3e81.png}
            \caption{Phân vùng Chuối bằng 1 Point Prompt}
        \end{figure}
    \end{columns}
\end{frame}

\section{Tổng kết}

\begin{frame}{Tổng kết Chương 11}
    \begin{itemize}
        \item \textbf{Image Segmentation} là bài toán tìm ranh giới chính xác đến từng pixel cho vật thể (Semantic, Instance, Panoptic).
        \item \textbf{U-Net \& Transposed Convolution:} Một mô hình kinh điển giải quyết phân vùng bằng kiến trúc Encoder (hạ phân giải để tìm Đặc trưng) và Decoder (tăng phân giải để Phục hồi không gian).
        \item \textbf{Segment Anything (SAM):} Đưa Computer Vision bước vào kỷ nguyên Foundation Model, cho phép phân vùng Zero-shot "Mọi thứ" chỉ bằng các gợi ý đơn giản (Point, Box, Text).
    \end{itemize}
\end{frame}

\end{document}
